% Canonical English version. Keep section/formula/citation structure identical
% to main_ru.tex.
\documentclass[10pt,twocolumn]{article}
\usepackage[margin=0.72in]{geometry}
\usepackage{amsmath,amssymb,booktabs,graphicx,microtype,url}
\usepackage[hidelinks]{hyperref}

\title{One Global Beam Across Many GPUs:\\
Exact High-Throughput Beam Search at Near-Billion Frontier Scale}
\author{Anonymous authors}
\date{}

\newcommand{\beam}{B}
\newcommand{\frontier}{\mathcal{F}}
\newcommand{\cand}{\mathcal{C}}

\begin{document}
\maketitle

\begin{abstract}
We present the first published high-throughput multi-GPU beam-search system
that physically shards one logical frontier across devices and jointly
constructs one next beam.  Candidate generation, scoring, 128-bit-key
de-duplication, routing, and frontier storage remain on the GPUs; only compact
control and ancestry metadata reach the host.  With no shard-local semantic
caps, the streamed procedure returns the same global reduced-key top-$B$ as a
monolithic implementation under the same score quantization, fixed-layout tie
order, and stated 128-bit identity model.  On eight A100 40GB GPUs, the
capacity result records a requested retained frontier of
$B_{\rm req}=770{,}883{,}178$; alignment gives
$B_{\rm eff}\geq B_{\rm req}$ and therefore at least
$18{,}501{,}196{,}272$ nominal pre-de-duplication children per saturated depth
for 24 generators.  That result contains no wall time or exact $B_{\rm eff}$
and is used only as capacity evidence.  Separately, a measured two-T4 run at
$B_{\rm eff}=82{,}837{,}504$ has $1{,}988{,}100{,}096$ nominal children per
saturated depth; depths 7--10 averaged $65.670$ s, or $30.274$ million logical
children/s end to end.  This throughput result covers two GPUs and does not
establish strong scaling beyond them.  The novelty is the full system-level
conjunction, not parallel, GPU, or distributed beam search in isolation.
\end{abstract}

\section{Introduction and Problem}
Beam search becomes a distributed-systems problem once the retained frontier
exceeds one GPU's memory.  Expansion, scoring, de-duplication, global
selection, and ancestry must be distributed without turning one search into
independent replicas or staging the frontier through host memory.  We ask
whether a GPU cluster can preserve the semantics of one monolithic beam while
scaling retained capacity with aggregate device memory.

Let $G=(V,E)$ be an implicit directed graph, $A_a(x)$ the child of state $x$
under generator $a\in\{1,\ldots,m\}$, and $\sigma(x,a)$ a deterministic
floating-point scoring oracle.  The mathematical contract and Streams~2--5
permit a learned or hand-engineered oracle provided Stream~1 emits deterministic
quantized keys; the shipped production Stream~1 backends and reported
experiments are neural.  The implementation ranks only the clipped and
quantized integer key $s(x,a)$ defined in Eq.~\eqref{eq:score}.  A
scalar state scorer is the special case $\sigma(x,a)=q(A_a(x))$, where
$q:V\rightarrow\mathbb{R}$; an output-$m$ scorer emits all action scores from
one parent evaluation.  At depth $d$, beam search computes
\begin{align}
  \cand_d &= \{(x,a,A_a(x)):x\in\frontier_{d-1},\ a\in[m]\},\nonumber\\
  U_d &= \operatorname{Reduce}_{z}
  (\cand_d;\,\operatorname{lexmin}(s,\mathrm{payload})),\\
  \frontier_d &= \operatorname{Top}_{\beam}
  \left(U_d;\,s_d^*,\tau_L\right),
  \label{eq:beam}
\end{align}
where $z$ is the 128-bit child-state key,
$s_d^*(z)=\min\{s(x,a):z(A_a(x))=z\}$, and $\tau_L$ is the
deterministic tie order induced by a fixed rank/shard layout $L$.  Thus
de-duplication keeps the best candidate for each key before the global cut.
For a saturated layer in which every retained parent is expanded by all $m$
generators, the nominal pre-de-duplication count is $m\beam_{\rm eff}$.  Since
runtime alignment gives $\beam_{\rm eff}\geq\beam_{\rm req}$, the measured
capacity request $\beam_{\rm req}=770{,}883{,}178$ with $m=24$ implies at least
$18{,}501{,}196{,}272$ nominal logical children in one depth.
Materializing its full 120-byte logical child states would require at least
$2.220$ TB before metadata; the design avoids that host staging and any
single-device global sort.

Our requirement is stronger than GPU-accelerated heuristic evaluation.  During
one depth, the candidate and frontier arrays in Eq.~\eqref{eq:beam} remain in
device memory.  The host receives only $O(p)$ control counts and compact
selected-candidate metadata, never the full state or frontier arrays.
Here ``one coherent search'' has an operational meaning: every GPU contributes
to Eq.~\eqref{eq:beam}, and the ranks collectively commit disjoint shards of
one globally defined next frontier.  It is not an ensemble of independent
beams or replicas.

\paragraph{Contributions.}
The system contributes (i) a GPU-resident hot-data path that distributes one
logical beam without frontier staging through host memory; (ii) exact
monolithic global reduced-key top-$B$ semantics under the same score
quantization, fixed-layout tie order, and stated 128-bit identity model, with
no shard-local semantic caps; and (iii) separate evidence for near-billion
retained capacity on eight GPUs and measured high-throughput execution on two
GPUs.  Together, these contributions make this the first published
high-throughput multi-GPU beam-search system with one physically sharded global
frontier and the stated exactness contract.  The implementation admits at most
128 ranks; throughput beyond two GPUs and execution beyond eight GPUs remain
unmeasured.

\section{Related Work and Gap}
Parallel beam search long predates GPUs: custom speech accelerators were studied
in 1986~\cite{anantharaman1986accelerator}.  Wijs and Lisser later distributed a
single global beam over CPUs~\cite{wijs2007distributed}, and Chong et al. placed
a data-parallel speech beam on one GPU~\cite{chong2008lvcsr}.  Thus, neither
parallel, distributed, nor GPU beam search is new by itself.  Modern CPU work
retains beams up to 20 million on a shared-memory machine, while MPI workers run
independent randomized beams~\cite{frohner2022parallel}; Kuroiwa and Beck also
study global and locally partitioned CPU variants~\cite{kuroiwa2024parallel}.

For GPU and heuristic combinatorial search, GA* establishes single-GPU A*
without a neural scorer~\cite{zhou2015gpuastar}, GPUAStar keeps neural A* on one
H200~\cite{soltani2026gpuastar}, and HDA* establishes hash-distributed CPU
best-first search~\cite{kishimoto2013hda}.  CayleyPy is the closest algorithmic
predecessor: its large-cube agents use one GPU and $B=2^{24}$
~\cite{chervov2025large,soibelman2025cayleypy}.  AlphaCube and DeepCubeA are
single-device beam search and batched weighted A*, respectively
~\cite{takano2023alphacube,agostinelli2019deepcubea}.

The closest adjacent systems result is QiankunNet-cuSCI, which generates
configurations, performs distributed global de-duplication, applies a neural
Transformer, and selects top-$K$ on 64 A100 GPUs with 91.06\% reported
strong-scaling efficiency~\cite{sun2026cunnqssci}.  It prevents a claim to the
generic expand--deduplicate--infer--select pipeline, but the work is not
reported as beam search.  It retains $K=32{,}000$ or $256{,}000$ selected SCI
configurations from roughly one billion raw candidates, and its global barriers
stage full candidate data through host memory.

Other billion-scale results measure different objects.  GPUexplore 3.0 explores
more than 3.15 billion reachable states on four GPUs but performs exhaustive
exploration without score-ranked bounded top-$B$~\cite{wijs2024gpuexplore}.  PathWeaver
pipelines approximate nearest-neighbor queries over a stored graph
~\cite{kim2025pathweaver}; multi-GPU speech and language decoders distribute
independent utterances, sentences, or model execution around tiny beams rather
than physically sharding one state frontier
~\cite{junczysdowmunt2016marian,braun2020gpuasr,chan2025trie}.

Table~\ref{tab:prior_art} applies one scale taxonomy: retained beam width is not
raw candidates, cumulative expansions, graph size, or a batch of independent
queries.  Prior systems establish important subsets of parallel beam search,
GPU execution, distributed search, or large-scale state exploration, but do
not report their full conjunction in one system.  This work is the first
published high-throughput multi-GPU beam-search system that combines one
physically sharded logical frontier, a GPU-resident candidate/frontier hot-data
path, global 128-bit-key de-duplication, no shard-local semantic caps, and exact
monolithic global reduced-key top-$B$ under the stated quantization, tie order,
and identity model at a near-billion retained-frontier request.  The priority
claim applies only to this full system-level conjunction.

\begin{table*}[t]
\centering
\footnotesize
\caption{Closest prior systems under a common scale taxonomy.  ``Other scale''
is deliberately not interpreted as retained beam width.}
\label{tab:prior_art}
\setlength{\tabcolsep}{2.5pt}
\begin{tabular}{@{}p{2.7cm}p{3.7cm}p{3.0cm}p{3.0cm}p{4.2cm}@{}}
\toprule
System & Platform and search object & Retained live beam & Other scale & Missing from the claimed conjunction \\
\midrule
Wijs--Lisser 2007 & distributed CPU; one global beam & 50,000 & 135.96M cumulative states & GPU data plane and device-resident scorer \\
Chong et al. 2008 & one 8800 GTX; speech active set & 20,000 cap & 1M-state model & multi-GPU state frontier \\
Frohner et al. 2022 & 46 CPU threads; one shared-memory beam & 20,000,000 & MPI: independent beams & GPU and distributed one-beam execution \\
CayleyPy 2025 & one GPU; one Cayley-graph beam & $2^{24}$ & independent agents & multi-GPU global frontier \\
GPUexplore 3.0 & four V100; exhaustive reachable set & not a beam & $>3.15$B states & ranked score and top-$B$ \\
PathWeaver 2025 & four A6000; independent ANNS queries & not reported & 50M-vector dataset & exact layered semantics and large retained beam \\
QiankunNet-cuSCI & 64 A100; selected SCI configurations & not a beam; $K=32{,}000/256{,}000$ & about 1B raw candidates & not reported as beam search; much smaller retained $K$; full-candidate host staging \\
This work & eight A100; one sharded global beam & request 770,883,178; effective $\geq$ request & $\geq18.501$B nominal raw children/depth & summary-backed capacity; time and effective width not recorded \\
\bottomrule
\end{tabular}
\end{table*}

\section{One Search Depth on Five GPU Streams}
The input is a global beam $\frontier_{d-1}$ distributed over $p$ GPUs.  If GPU
$r$ owns $N_r$ parents and the graph has $m$ moves, it generates $mN_r$
logical candidates.  For $y=A_a(x)$, a candidate is represented before
materialization by
\begin{equation}
 c=(z(y),s(x,a),i_{\mathrm{parent}},a,r_{\mathrm{src}}),
 \label{eq:candidate}
\end{equation}
where $z(y)\in\{0,1\}^{128}$ is the child-state hash, $s(x,a)$ is a
quantized candidate score, and $r_{\mathrm{src}}$ identifies the GPU that stores
the parent.  The child state itself is reconstructed only for goal checking or
after the final selection.

\paragraph{Streams 1--2: score and identity.}
Parents are divided into microbatches and placed in reusable ring slots.
Streams~1 and~2 consume the same slot concurrently.  Stream~1 evaluates the
chosen scoring oracle and writes only
\begin{equation}
 s(x,a)=\operatorname{rint}_{\mathrm{RN-even}}
 \left(S\,\operatorname{clip}(\sigma(x,a),0,q_{\max})\right)
 \in\mathbb{N},
 \label{eq:score}
\end{equation}
not the floating-point output.  Here RN-even is the GPU path's default
round-to-nearest, ties-to-even mode.  The scalar path evaluates child states, whereas
the output-$m$ path obtains all $m$ action scores from each parent.  Stream~2
applies every move in registers, checks the goal, and writes the 128-bit Zobrist
hash.  Thus score evaluation and state manipulation overlap without storing the
$mN_r$ child states.

\paragraph{Stream 3: early reduction and ownership.}
When a group of ring slots has both scores and hashes, Stream~3 removes scores
above the current threshold $T$, compacts the survivors, sorts them by $z$, and
reduces equal hashes.  For each hash it retains the lexicographically smallest
pair $(s,\mathrm{payload})$, where the payload recovers the original parent and
move.  It then assigns the unique candidate to
\begin{equation}
 o(x)=\operatorname{mix}(z(x))\bmod p.
 \label{eq:owner}
\end{equation}
Local candidates enter the local collector; the rest are grouped by owner.
Performing this first reduction before communication lowers traffic, but does
not replace the later global deduplication.

\paragraph{Stream 5: GPU-to-GPU exchange.}
Stream~5 exchanges the owner-grouped metadata directly between GPUs using
NCCL.  The owner rule implies
\begin{equation}
 z(x)=z(y)\Longrightarrow o(x)=o(y),
 \label{eq:owner_invariant}
\end{equation}
so duplicates produced by different source GPUs eventually meet at one owner.
Stream~5 performs communication only; received candidates re-enter the same
collector as local candidates.

\paragraph{Stream 4: persistent survivor set.}
The collector partitions owned candidates into logical shards.  Each shard has
two resident buffers: Stream~3 appends to one while Stream~4 cleans the other.
A Stream~4 job merges the clean and newly appended regions, applies its snapshot
of $T$, compacts, sorts by $z$, and retains the best metadata for each hash.
There is no per-shard top-$k$ and no semantic shard quota: a shard may retain
every candidate below $T$.  Consequently sharding does not approximate the
score cut; only the equal-score boundary order is defined by the fixed layout.

\paragraph{Classical top-$\beam$ versus monotone thresholding.}
Let $C$ be the complete candidate multiset, write $s(c)$ for the quantized
integer score in Eq.~\eqref{eq:score}, and let
$U=\operatorname{Reduce}_{z}(C;\operatorname{lexmin}(s,\mathrm{payload}))$ be
its one-entry-per-key reduction, carrying
$s^*(z)=\min_{c:z(c)=z}s(c)$ and the associated best payload.  For fixed layout
$L$, rank keys by $\kappa_L(z)=(s^*(z),\tau_L(z))$, where smaller score is
better and $\tau_L$ is the implementation's deterministic
rank-then-physical-order rule for equal scores.  A classical reference
materializes $U$, sorts it by $\kappa_L$, and returns, for
$b=\min(\beam,|U|)$,
\begin{equation}
 \operatorname{Top}_{\beam}(U)=
 \{z\in U:\kappa_L(z)\leq \kappa_{L,(b)}(U)\}.
 \label{eq:classical_topb}
\end{equation}
Our implementation never materializes $U$.  After update $j$, let
$C_j\subseteq C_{j+1}\subseteq C$ be the candidates processed so far and
$U_j=\operatorname{Reduce}_{z}(C_j)$.  The key domain of $U_j$ can only grow,
while an existing key score $s_j^*(z)$ can only decrease.  Once
$|\operatorname{dom}U_j|\geq\beam$, histograms of the current reduced scores
define
\begin{align}
 \widehat T_j&=\min\left\{t:\sum_{u=0}^{t}H_j(u)\geq\beam\right\},\nonumber\\
 T_j&=\min\{T_{j-1},\widehat T_j\},\qquad T_0=+\infty.
 \label{eq:online_threshold}
\end{align}
While $|\operatorname{dom}U_j|<\beam$, $\widehat T_j=+\infty$.
Streams~3 and~4 may discard only candidates with $s(c)>T_j$; the score boundary
$s(c)=T_j$ is retained because its final membership depends on $\tau_L$.

\paragraph{Exactness.}
Processing another candidate either introduces a key or lowers an existing
key's best score, so the $\beam$-th smallest reduced score cannot increase;
hence $\widehat T_{j+1}\leq\widehat T_j$ and $T_{j+1}\leq T_j$.  If candidate
$c$ is discarded at update $j$, then $s(c)>T_j$ and $U_j$ already contains at
least $\beam$ distinct keys with current best scores at most $T_j$.  Those keys
remain distinct and their final best scores can only decrease.  If $z(c)$ is
among them, $c$ is worse than its stored representative; otherwise at least
$\beam$ other keys are strictly better.  Thus $c$ cannot determine a selected
final minimum.  Conversely, a candidate attaining the final minimum of a
global top-$\beam$ key cannot be discarded, because that would require
$\beam$ distinct keys already carrying strictly smaller scores.  Thresholding
therefore preserves exactly the classical reduced-key score top-$\beam$ under
the same quantization for fixed $L$, subject to the same fixed-layout tie rule
at the boundary.

\paragraph{Identity and exactness scope.}
The proof is exact for the score quantization in Eq.~\eqref{eq:score}, the set
keyed by the 128-bit Zobrist identifier, and the complete fixed-layout order
$\kappa_L$.  The candidate data plane does not carry all 120
logical state bytes, so equal 128-bit keys are treated as equal states rather
than resolved by a full-state comparison.  Selection is therefore
mathematically exact under this stated key model; the residual risk is a
128-bit hash collision, not an approximation introduced by sharding,
thresholding, or top-$\beam$.  Changing the GPU or shard layout can change
$\tau_L$ and therefore which equal-score states occupy the boundary; the
current implementation does not claim partition-invariant selected IDs.

\paragraph{Final cut and materialization.}
After all five streams are drained, the ranks combine the final histograms.  If
$|U|\geq\beam$, they choose the smallest $T^*$ such that
\begin{equation}
 \sum_{u=0}^{T^*}H(u)\geq\beam.
 \label{eq:threshold}
\end{equation}
All states below $T^*$ survive, and a deterministic tie-break cuts its boundary
to exactly $\beam$ states.  If $|U|<\beam$, all available keys survive instead.
Thus the next frontier always contains exactly $\min(\beam,|U|)$ states.  The
selected metadata are load-balanced across GPUs.
Only then does each source GPU reconstruct the requested children and send
them directly to their target GPUs.  The CPU receives $O(p)$ count vectors plus
compact selected-candidate metadata used to reconstruct ancestry, but no full
state or frontier arrays.

The five streams are readiness-driven rather than five serial phases.  Streams
1--2 fill ring slots almost continuously while Streams~3 and~4 reduce earlier
batches and Stream~5 exchanges remote survivors.  The mathematical invariant is
that every hash surviving the final cut is reduced at its unique owner before
Eq.~\eqref{eq:threshold}; hence, for a fixed layout, the distributed execution
implements the same global reduced-key score top-$\beam$ as its monolithic
reference under the same score quantization, without materializing the raw
candidate layer.

\section{Experiments}
\paragraph{Protocol.}
We evaluate computational cost rather than solution quality and report a
measurement only with its hardware, workload, model contract, and scope.
Production time includes generation, de-duplication, communication, selection,
and history; it is not interchangeable with an isolated Stream~1 rate.  Empty
cells mean ``not measured'', never an inferred value.  Table~\ref{tab:platforms}
summarizes the strongest locally verified measurements.

\begin{table*}[t]
\centering
\footnotesize
\caption{Measured performance by platform.  E2E is the production critical
path; micro is an isolated kernel pipeline; capacity has no timing claim.
n/d: not derived; n/a: not applicable; n/r: not recorded.}
\label{tab:platforms}
\setlength{\tabcolsep}{3pt}
\begin{tabular}{@{}p{2.0cm}p{2.4cm}p{3.8cm}p{3.2cm}p{1.8cm}p{1.8cm}@{}}
\toprule
Platform & Workload & Beam $B$ & Measurement & s/depth & Raw child rate \\
\midrule
RTX 3070 8GB & puzzle 0, depth 20 & $2^{22}$ & E2E complete; $60.623$ s total & 3.031 avg. & n/d \\
RTX 3070 8GB & Stream 1, $B_\mu=4096$, $c=3$ & --- & micro complete & n/a & $48.142$ M/s \\
Kaggle $2\times$T4 & puzzle 20, radius 4, limit 72 & $B_{\rm req}=82{,}615{,}524$; $B_{\rm eff}=82{,}837{,}504$ & solved; $350.402$ s; CSV+summary & 65.670 saturated & $30.274$ M/s \\
$8\times$A100 40GB & puzzles 900/950/1000, depth 8 & $B_{\rm req}=770{,}883{,}178$; $B_{\rm eff}$ not recorded & stable; $39{,}745$ MiB/GPU; summary & n/r & n/r \\
\bottomrule
\end{tabular}
\end{table*}

In Table~\ref{tab:platforms}, the RTX value is total E2E wall divided by 20
completed depths, whereas the Kaggle value is the summary-reported mean of
saturated depths 7--10.  The T4 CSV records both requested and aligned effective
widths; $30.274$ M/s is the explicit derivation
$24\times82{,}837{,}504/65.670$ and counts logical children before
de-duplication.  We do not derive an RTX E2E rate because its early depths are
not saturated.  Kernel microbenchmarks and the A100 capacity-only summary do
not define an E2E depth time and are therefore marked not applicable or not
recorded.

\paragraph{Local RTX 3070.}
On the laptop GPU, the folded-input half2 kernel raised the best isolated
Stream~1 rate from approximately $24.7$ to $48.142$ million candidates/s.  In
Nsight, the changed kernel fell from $9.678$ to $3.232$ s over 1,202 launches
($2.99\times$).  More importantly, the same production depth-20 workload fell
from $115.334$ to $60.623$ s, a $1.90\times$ end-to-end improvement.  A legacy
PyTorch run at depth 100 exhausted memory at depth 33 after $329.067$ s; we do
not count this incomplete run as a speedup baseline.

\paragraph{Kaggle $2\times$T4.}
The two ranks jointly maintain one global beam.  For puzzle 20, requested
$B_{\rm req}=82{,}615{,}524$ was aligned to
$B_{\rm eff}=82{,}837{,}504$ and the radius-4-neighborhood run solved in
$350.402$ s with 128 shards and $B_\mu=4096$.  At $B=2^{26}$, increasing the shard count from 8 to 64 reduced
the mean time of saturated depths 7--9 from $69.365$ to $51.732$ s
($1.34\times$) while retaining exact global selection.  The best isolated
Stream~1 point was $40.60$ M candidates/s; it is not used as an E2E rate.

\paragraph{$8\times$A100 capacity and scale-out limit.}
The exact-scale autotune summary records stable depth-8 execution on
three Megaminx instances for requested
$B_{\rm req}=770{,}883{,}178$, peaking at $39{,}745$ MiB per GPU.  Runtime
alignment makes $B_{\rm eff}\geq B_{\rm req}$, but the exact effective value,
raw rank logs, runtime configuration, commit, checkpoint hash, and wall time
were not preserved.  The summary therefore supports an eight-device capacity
result and a nominal saturated scale of at least $18.501$ billion logical child
candidates, not a near-billion speed result.  Performance is measured only
through two GPUs and capacity through eight; static exchange tables currently
admit at most 128 ranks, but execution beyond eight GPUs is unmeasured.

\section{Comparison with CayleyPy and Megaminx Result}
CayleyPy establishes empirically that wider beams improve success and solution
quality, but its published examples remain single-device.  Its controlled
$S_{20}$ experiment increases $W$ from $2^{10}$ (no solution) to $2^{20}$
(the ideal length 190), while the large-cube paper uses $W=2^{24}$ so that an
agent fits on one GPU~\cite{soibelman2025cayleypy,chervov2025large}.  Our
requested $770{,}883{,}178$ configuration is $735.2\times$ wider than $2^{20}$ and
$45.95\times$ wider than $2^{24}$.  It is also $38.54\times$ wider than the
largest peer-reviewed retained beam located in our audit (20 million) and
$30.84\times$ wider than the 25-million public HATETRIS
implementation~\cite{hatetrispublic}.  These
ratios compare systems scale, not
solution difficulty across different graphs.

\begin{table}[t]
\centering
\small
\caption{Published CayleyPy beam-width trend (Table 8 of
\cite{soibelman2025cayleypy}).}
\label{tab:cayleypy_width}
\begin{tabular}{@{}rrrrrr@{}}
\toprule
$W$ & $2^{10}$ & $2^{12}$ & $2^{14}$ & $2^{16}$ & $2^{20}$ \\
\midrule
Length & --- & 252 & 202 & 194 & 190 \\
\bottomrule
\end{tabular}
\end{table}

We additionally measured Pilgrim and our search implementation on the same
Cube4 state, checkpoint, single GPU, beam, and ten layers.  The scorer is a
3,383,064-parameter piece Transformer with 96 six-class positions,
$d_{\rm model}=256$, four layers, eight heads, feed-forward width 1024, and 24
Q outputs; inference uses FP16 and microbatches of 384.  Table~\ref{tab:cube4_pair}
reports the tenth layer, where every saturated frontier expands $24B$ pairs.

\begin{table}[t]
\centering
\small
\setlength{\tabcolsep}{2pt}
\caption{Matched Cube4 depth-10 computation at $B=4{,}194{,}304$.}
\label{tab:cube4_pair}
\resizebox{\columnwidth}{!}{%
\begin{tabular}{@{}llrrr@{}}
\toprule
GPU & Implementation & s/depth (depth 10) & Pairs/s & Peak VRAM \\
\midrule
T4 & Pilgrim & 198.153 & 508,007 & 2.90 GiB reserved \\
T4 & Ours & 181.282 & 555,293 & 1,949 MiB used \\
P100 & Pilgrim & 314.487 & 320,087 & 2.90 GiB reserved \\
P100 & Ours & 301.955 & 333,378 & 1,949 MiB used \\
\bottomrule
\end{tabular}%
}
\end{table}

\subsection{Native approach inventory and matched model families on one T4}
Table~\ref{tab:four_t4} is an inventory, not a cross-row speed ranking.
AlphaCube and DeepCubeA retain their released Cube3 models. Megaminx results
form two disjoint matched families: scalar \texttt{output\_dim=1} (Pilgrim
Searcher, native CayleyPy, MultiGPUBeamSearch) and native
\texttt{output\_dim=24} (Pilgrim QSearcher, MultiGPUBeamSearch). Ratios require
the identical checkpoint, architecture, head, state, generators, depth, width,
dtype, and hardware. CayleyPy cannot consume the 24-output head and remains
incompatible rather than transformed.
The rate column contains only native counted nodes divided by solver wall time;
it is never replaced by $B/t$.

\begin{table*}[t]
\centering
\small
\setlength{\tabcolsep}{3pt}
\caption{Exact neural architectures used by the reported tests. Parameter
counts refer to the trainable source model before batch-normalization folding.}
\label{tab:model_architectures}
\resizebox{\textwidth}{!}{%
\begin{tabular}{@{}lrlp{10.8cm}@{}}
\toprule
Model family & Parameters & Used by & Ordered block sequence \\
\midrule
Cube3 AlphaCube 0.1.6 \texttt{large} & 118,939,666 & AlphaCube & one-hot 324 $\to$ $8\times$[Linear 4096, ReLU, BN] $\to$ Linear 18 \\
Cube3 DeepCubeA official \texttt{cube3} & 14,663,001 & DeepCubeA & one-hot 324 $\to$ Linear 5000, BN, ReLU $\to$ Linear 1000, BN, ReLU $\to$ $4\times$[Linear 1000, BN, ReLU, Linear 1000, BN, residual add, ReLU] $\to$ Linear 1 \\
Megaminx scalar MLP & 34,766,849 & Pilgrim Searcher; CayleyPy; MultiGPU & position--class $120\times120$ $\to$ EmbeddingBag-equivalent Linear 2048, BN, ReLU $\to$ Linear 512, BN, ReLU $\to$ $8\times$[Linear 512, BN, ReLU, Linear 512, BN, residual add, ReLU] $\to$ Linear 1 \\
Megaminx QMLP2RB & 23,978,008 & Pilgrim QSearcher; MultiGPU & position--class $120\times120$ $\to$ EmbeddingBag-equivalent Linear 1536, BN, ReLU $\to$ Linear 512, BN, ReLU $\to$ $2\times$[Linear 512, BN, ReLU, Linear 512, BN, residual add, ReLU] $\to$ Linear 24 \\
Cube4 Q Transformer & 3,383,064 & matched Cube4 only & 96 tokens / 6 classes $\to$ projection + CLS $\to$ $4\times$[pre-LN, 8-head attention ($d=256$), residual, pre-LN, FFN 256--1024--256 ReLU, residual] $\to$ LN $\to$ Linear 24 \\
\bottomrule
\end{tabular}%
}
\end{table*}

For the matched Megaminx rows, all listed implementations load the same bytes:
the scalar checkpoint has SHA-256 \texttt{7f5071e6...a539b759}, while the
output-24 checkpoint has SHA-256 \texttt{e7bda332...5d2ca670}. The latter is
not derived from the scalar model. The Cube4 Transformer is a separate matched
comparison and is never mixed into the Megaminx capacity sweep.

\begin{table*}[t]
\centering
\small
\setlength{\tabcolsep}{3pt}
\caption{Native search inventory on one visible T4. Rows with different model
contracts are not performance comparisons; ``Max complete'' is measured.}
\label{tab:four_t4}
\resizebox{\textwidth}{!}{%
\begin{tabular}{@{}lllrrrl@{}}
\toprule
Approach & Workload / native scale & Max complete & Wall (s) & s/native step & Nodes/s & Boundary/status \\
\midrule
AlphaCube 0.1.6 & Cube3 \texttt{large}, beam & 1,572,864 & 461.488 & 27.146/depth & 42,387 & 1,703,936 host SIGKILL \\
DeepCubeA & Cube3 A*, 55 iterations & n/a & 146.474 & 2.663/iteration & 41,880 & complete, length 21 \\
Pilgrim Searcher & Megaminx scalar beam, depth 8 & 5,898,240 & 1,462.557 & 182.820/depth & not measured & 6,291,456 OOM \\
Pilgrim QSearcher & Megaminx output-24 beam, depth 8 & 20,971,520 & 197.358 & 24.670/depth & not measured & 23,068,672 OOM \\
native CayleyPy & Megaminx scalar beam, depth 8 & 20,971,520 & 4,258.602 & 532.325/depth & not measured & broad: $2^{24}$ complete, $2^{25}$ OOM \\
MultiGPUBeamSearch & Megaminx scalar beam, depth 8 & 29,360,128 & 2,817.158 & 352.145/depth & not measured & 31,457,280 OOM \\
MultiGPUBeamSearch & Megaminx output-24 beam, depth 8 & 29,360,128 & 116.904 & 14.613/depth & not measured & 33,554,432 OOM \\
\bottomrule
\end{tabular}%
}
\end{table*}

\begin{table}[t]
\centering
\small
\caption{AlphaCube 1xT4 completed sweep.}
\label{tab:alpha_curve}
\begin{tabular}{@{}rrrrr@{}}
\toprule
$B$ & Wall (s) & s/depth & Nodes & Nodes/s \\
\midrule
1,024 & 1.034 & 0.052 & 18,720 & 18,108 \\
4,096 & 1.199 & 0.060 & 73,484 & 61,309 \\
16,384 & 5.014 & 0.251 & 282,380 & 56,321 \\
65,536 & 18.543 & 1.091 & 903,448 & 48,722 \\
131,072 & 37.818 & 2.224 & 1,755,416 & 46,418 \\
262,144 & 79.166 & 4.657 & 3,459,352 & 43,698 \\
524,288 & 170.671 & 10.039 & 6,867,224 & 40,237 \\
786,432 & 228.938 & 13.467 & 10,124,136 & 44,222 \\
1,048,576 & 303.870 & 17.875 & 13,269,864 & 43,670 \\
1,310,720 & 382.158 & 22.480 & 16,415,592 & 42,955 \\
1,572,864 & 461.488 & 27.146 & 19,561,320 & 42,387 \\
\bottomrule
\end{tabular}
\end{table}

\begin{table*}[t]
\centering
\small
\setlength{\tabcolsep}{4pt}
\caption{Completed native CayleyPy 1xT4 powers-of-two sweep. Memory is the
PyTorch CUDA allocator peak; the final row is the first broad-grid failure.}
\label{tab:cayleypy_curve}
\begin{tabular}{@{}rrrrrr@{}}
\toprule
$B$ & Wall (s) & s/depth & Allocated (MiB) & Reserved (MiB) & Status \\
\midrule
1,024 & 0.930 & 0.116 & 85.9 & 92.0 & complete \\
2,048 & 1.237 & 0.155 & 96.3 & 112.0 & complete \\
4,096 & 2.189 & 0.274 & 117.3 & 128.0 & complete \\
8,192 & 3.828 & 0.478 & 158.6 & 230.0 & complete \\
16,384 & 6.776 & 0.847 & 241.5 & 292.0 & complete \\
32,768 & 13.722 & 1.715 & 407.5 & 510.0 & complete \\
65,536 & 31.399 & 3.925 & 739.3 & 858.0 & complete \\
131,072 & 63.664 & 7.958 & 772.2 & 1,180.0 & complete \\
262,144 & 114.551 & 14.319 & 834.8 & 1,168.0 & complete \\
524,288 & 219.991 & 27.499 & 962.1 & 2,156.0 & complete \\
1,048,576 & 430.789 & 53.849 & 1,216.7 & 3,234.0 & complete \\
2,097,152 & 745.341 & 93.168 & 1,726.0 & 2,740.0 & complete \\
4,194,304 & 1,320.000 & 165.000 & 2,745.2 & 4,492.0 & complete \\
8,388,608 & 2,489.356 & 311.169 & 5,155.3 & 6,462.0 & complete \\
16,777,216 & 4,741.581 & 592.698 & 10,233.7 & 14,652.0 & complete \\
33,554,432 & --- & --- & 12,922.2 & 14,484.0 & CUDA OOM \\
\bottomrule
\end{tabular}
\end{table*}

For AlphaCube, s/depth is solver wall divided by the completed search depth
reported by the returned solution (20 for the first three rows and 17 for the
remaining rows).  For both Megaminx implementations it is wall divided by the
fixed eight completed depths.  These averages include startup and early
unsaturated layers; they complement, rather than replace, nodes/s.

Pilgrim's widest scalar completion, $B=5{,}898{,}240$, took 1,462.557 s
(182.820 s/depth) with a 14,433 MiB device peak; $B=6{,}291{,}456$ CUDA-OOMed.
For output-24, $B=20{,}971{,}520$ completed in 197.358 s
(24.670 s/depth) at 12,333 MiB, whereas $B=23{,}068{,}672$ CUDA-OOMed. The
measured brackets are therefore 5,898,240--6,291,456 and
20,971,520--23,068,672, respectively.

The final collaborator-run MultiGPU pass and scalar-only v27 confirmation now
give measured complete/OOM pairs for both heads. Scalar completed at
29,360,128 and CUDA-OOMed at 31,457,280; native output-24 completed through
29,360,128 and CUDA-OOMed at 33,554,432.

\begin{table}[t]
\centering
\small
\caption{Final MultiGPU 1xT4 scalar boundary points.}
\label{tab:multigpu_v11_partial}
\resizebox{\columnwidth}{!}{%
\begin{tabular}{@{}rrrrl@{}}
\toprule
$B$ & Wall (s) & s/depth & Peak device (MiB) & Status \\
\midrule
25,165,824 & 3,027.673 & 378.459 & 12,267 & complete \\
27,262,976 & 3,080.206 & 385.026 & 13,241 & complete \\
29,360,128 & 2,817.158 & 352.145 & 14,215 & complete \\
31,457,280 & --- & --- & 14,909 & CUDA OOM \\
\bottomrule
\end{tabular}
}
\end{table}

The scalar OOM recurred for every row-budget profile from 49,152 through 24.
The common primary metric is total wall divided by eight, 352.145 s/depth;
the auxiliary final-depth time was 1,644.13 s.

\begin{table}[t]
\centering
\small
\caption{Final MultiGPU 1xT4 native output-24 sweep.}
\label{tab:multigpu_output24}
\resizebox{\columnwidth}{!}{%
\begin{tabular}{@{}rrrrl@{}}
\toprule
$B$ & Wall (s) & s/depth & Peak device (MiB) & Status \\
\midrule
1,048,576 & 8.296 & 1.037 & 773 & complete \\
4,194,304 & 22.735 & 2.842 & 2,253 & complete \\
8,388,608 & 41.091 & 5.136 & 4,159 & complete \\
16,777,216 & 77.460 & 9.683 & 7,983 & complete \\
20,971,520 & 96.124 & 12.015 & 9,893 & complete \\
25,165,824 & 111.130 & 13.891 & 11,799 & complete \\
29,360,128 & 116.904 & 14.613 & 13,717 & complete \\
33,554,432 & 1.563 & --- & 14,909 & CUDA OOM \\
\bottomrule
\end{tabular}
}
\end{table}

The output-24 OOM recurred while the runtime descended its row-budget profiles
from 2,048 to 1. Its 1.563-s field is the final failed attempt, not a completed
depth time.

\begin{table}[t]
\centering
\scriptsize
\setlength{\tabcolsep}{3pt}
\caption{Matched 1xT4 Megaminx scalar sweep. Neither implementation exposes a
valid node counter for these points.}
\label{tab:megaminx_curve}
\begin{tabular}{@{}rrrrr@{}}
\toprule
& \multicolumn{2}{c}{MultiGPU} & \multicolumn{2}{c}{CayleyPy} \\
\cmidrule(lr){2-3}\cmidrule(l){4-5}
$B$ & Wall (s) & s/depth & Wall (s) & s/depth \\
\midrule
65,536     & 12.537    & 1.567   & 31.399    & 3.925 \\
262,144    & 39.935    & 4.992   & 114.551   & 14.319 \\
1,048,576  & 152.604   & 19.076  & 430.789   & 53.849 \\
4,194,304  & 476.007   & 59.501  & 1,320.000 & 165.000 \\
8,388,608  & 892.948   & 111.619 & 2,489.356 & 311.169 \\
16,777,216 & 1,728.460 & 216.058 & 4,741.581 & 592.698 \\
20,971,520 & 2,179.990 & 272.499 & 4,258.602 & 532.325 \\
25,165,824 & 2,416.150 & 302.019 & ---       & --- \\
\bottomrule
\end{tabular}
\end{table}

AlphaCube materializes all $18B$ Cube3 children in host NumPy arrays before
sorting; consequently its first failure is a host-memory kill even though GPU
reservation is only 4.685~GiB.  Its measured rate settles near 40--44 thousand
nodes/s at the widest completed points.  DeepCubeA has no persistent beam: it
is weighted A* and its native scale is search/NN batch 10,000/10,000, so
inventing a frontier maximum for it would be incorrect.  Its completed T4 run
performed 55 search iterations in 146.474 s, or 2.663 s per native A* iteration;
this is not a beam-depth measurement.  The same run expanded 6,134,352 nodes
at 41,880 nodes/s and returned a length-21 solution.

For the matched scalar Megaminx pair, CayleyPy commit \texttt{5b2b53f} uses
iterated search with \texttt{history\_depth=0}. The broad protocol completed
every power of two from $2^{10}$ through $2^{24}$ and CUDA-OOMed at $2^{25}$;
the memory-trend estimate equalled $2^{24}$, so no redundant confirmation was
scheduled. Earlier targeted evidence with the same contract retains the finer
20,971,520-complete versus 21,037,056-OOM bracket, but is not presented as part
of the broad-grid policy. Its internal expansion path
bypassed the public counter hook, so only wall time and memory are valid.  The
MultiGPU likewise retained end-to-end wall but no valid node counter. The
final scalar pass measured 29,360,128 complete at 14,215 MiB versus
31,457,280 CUDA OOM at 14,909 MiB. The matched output-24 pass measured
29,360,128 complete versus 33,554,432 CUDA OOM. Superseded static
profile failures and the cancelled v10/v11 runs remain excluded.
At the latest broad-grid matched width, CayleyPy took 4,741.581 s and the
historical MultiGPU run took 1,728.460 s at $B=16{,}777{,}216$ (2.74x lower).
At the older targeted $B=20{,}971{,}520$ point, MultiGPU wall is 1.95x lower.
These are matched
end-to-end wall ratios, not inferred node-throughput ratios.

The scalar ratios above use one identical scalar model contract and must not be
compared with the faster Pilgrim QSearcher row. The 24-output model evaluates
all 24 actions per parent and has a different architecture; scalar-versus-Q
timing is therefore a model-contract difference, not an algorithmic speedup.
The separate Cube4 comparison uses the identical 24-output Q checkpoint for
Pilgrim and MultiGPUBeamSearch. No output-24 CayleyPy value is inferred.

Thus, at the matched Pilgrim point our
tenth-layer throughput is $1.093\times$ on T4
and $1.042\times$ on P100.  PyTorch allocator reservation and device-wide
\texttt{nvidia-smi} usage are different memory metrics, so the table does not
claim a memory ratio; it establishes that both implementations complete 4M.

For the Megaminx superflip instance, the located public chronology improves
from 55 through 54 and 53 to 51 moves; the immediately preceding public value
was therefore 53, not 55~\cite{megaminxsuperflip53,megaminxsuperflip51}.  The
51-step path uses the same 24-generator signed 72-degree face-turn metric and
has public wall time $22{,}600.2$ s.  Independent CPU replay of the recorded
submission reaches the exact central state with zero differences; its path
and file SHA-256 values are recorded in the accompanying evidence audit.
The cited public chronology contains no shorter same-metric path, but no
formal registry or optimality proof exists.  We therefore call it the shortest
public same-metric solution located, not an unqualified world record.  Its
missing run manifest is kept separate from the summary-backed
$770{,}883{,}178$-requested-state depth-8 capacity report.

\section{Code and Artifacts}
Paper sources, the evidence index, and machine-readable tables are public at
\url{https://github.com/TryDotAtwo/MultiGPUBeamSearchPaper}.  Release v1.0.0
contains both PDFs and the matched checkpoints:
\url{https://github.com/TryDotAtwo/MultiGPUBeamSearchPaper/releases/tag/v1.0.0}.
Solver source and its GPU-ready Megaminx release are public at
\url{https://github.com/TryDotAtwo/MultiGPUBeamSearch} and
\url{https://github.com/TryDotAtwo/MultiGPUBeamSearch/releases/tag/megaminx-native-cluster-latest}.
The public self-contained one-T4 reproduction dataset is
\url{https://www.kaggle.com/datasets/trydotatwo/gpu-beam-search-paper-1xt4-reviewer-notebooks}.
Four importable notebook sources for the Pilgrim and MultiGPUBeamSearch broad
sweeps and boundary confirmations, together with CayleyPy, AlphaCube, and
DeepCubeA runners, are versioned in the paper repository.  They consume the
public dataset through Kaggle's normal Import Notebook and Save Version
workflow and do not depend on the visibility of an owner's saved kernel page.
The eight-A100 record supplies capacity and memory evidence but no wall time or
exact effective width; the two-T4 record separately supplies measured
throughput.  No missing timing, capacity, or scaling metric is inferred from a
different run.

\section{Conclusion}
This work establishes one GPU-resident distributed beam search rather than a
collection of local surrogates.  For a fixed layout, its streamed
expand--de-duplicate--select path, with no shard-local semantic caps, is
mathematically equivalent to monolithic global reduced-key top-$B$ under the
same score quantization, tie order, and stated 128-bit identity model.  The
eight-A100 result establishes a requested frontier of $770{,}883{,}178$
retained states and at least $18.501$ billion nominal pre-de-duplication
children per saturated depth; it does not provide wall time or exact effective
width.  The separate two-T4 result measures $30.274$ million logical children/s
end to end, so high-throughput execution is demonstrated through two GPUs while
stronger scaling remains unmeasured.  The project also provides a replay-valid
51-move Megaminx-superflip path, the shortest public same-metric solution
identified in the supporting evidence.  This is the first published
high-throughput multi-GPU beam-search system with one physically sharded global
frontier and the stated exactness contract at near-billion retained-frontier
scale.

\begingroup
\footnotesize
\enlargethispage{4\baselineskip}
\bibliographystyle{abbrv}
\bibliography{references}
\endgroup
\end{document}
